% Separated-block resultants: rigorous G3 results and the exact G4 obstruction.
% This is a fragment intended to be read with proof.tex and nv_theorem.tex.

\subsection{Separated root strips and centered resultant structure}

Retain the notation
\[
 P(n)=(2n+1)(17n^2+17n+5),\qquad
 N_{h+1}(x)=P(x+h)N_h(x)-(x+h)^6N_{h-1}(x),
\]
with $N_0=0$ and $N_1=1$, and, for $d,r\geq2$, put
\[
 \mathcal R_{d,r}
 =\operatorname {Res}_x\bigl(N_d(x),N_r(x+d)\bigr).
\]
We use the standard resultant convention.  The first result proves, in
particular, that every integer $\mathcal R_{d,r}$ is nonzero.

We will use only the absolute adjacent-resultant formula from
Remark~\ref{rem:adj-res}.  With the standard signed convention, its
one-step recursion for $h\geq2$ has the correction
\[
 \operatorname {Res}(N_h,N_{h+1})
 =(-1)^{h-1}N_h(-h)^6\operatorname {Res}(N_{h-1},N_h);
\]
the factor $(-1)^{h-1}$ comes from the scalar minus sign in the recurrence.

\begin{proposition}[Separated root strips]\label{prop:meso-root-strip}
For $d\geq2$, every complex root $\alpha$ of $N_d$ satisfies
\begin{equation}\label{eq:meso-root-strip}
 -d<\operatorname {Re}\alpha<-1.
\end{equation}
Consequently, for all $d,r\geq2$,
\[
 \mathcal R_{d,r}>0.
\]
\end{proposition}

\begin{proof}
Let $\mathsf A_d(x)$ be the $(d-1)\times(d-1)$ symmetric tridiagonal matrix
with
\[
 (\mathsf A_d)_{ii}=P(x+i),\qquad
 (\mathsf A_d)_{i,i+1}=(\mathsf A_d)_{i+1,i}=(x+i+1)^3.
\]
Expansion at the last row gives the defining continuant recurrence, so
$\det\mathsf A_d(x)=N_d(x)$.

Suppose $\operatorname {Re}x\geq-1$ and set
$z_i=x+i+\tfrac12$.  Then $|z_i|\geq\tfrac12$ and
\[
 P(x+i)=34z_i^3+\frac32z_i.
\]
It follows that
\[
 |P(x+i)|
 \geq |z_i|\bigl(34|z_i|^2-\tfrac32\bigr)
 \geq28|z_i|^3.
\]
The sum of the absolute values of the two possible off-diagonal entries
in row $i$ is at most
\[
 |z_i-\tfrac12|^3+|z_i+\tfrac12|^3
 \leq2(|z_i|+\tfrac12)^3
 \leq16|z_i|^3.
\]
Thus $\mathsf A_d(x)$ is strictly row-diagonally dominant and is invertible.
There is no root with $\operatorname {Re}x\geq-1$.

If $J$ is the reversal matrix, the definitions and
$P(-z-1)=-P(z)$ give
\[
 \mathsf A_d(-x-d-1)=-J\mathsf A_d(x)J.
\]
Taking determinants proves the reflection identity
\[
 N_d(-x-d-1)=(-1)^{d-1}N_d(x)
\]
then excludes $\operatorname {Re}x\leq-d$, proving
\eqref{eq:meso-root-strip}.  The roots of $N_r(x+d)$ lie in the
disjoint strip
\[
 -d-r<\operatorname {Re}x<-d-1,
\]
so the two polynomials have no common complex root and their resultant
is nonzero.

Finally, the leading coefficients $\ell_d,\ell_r$ are positive.  In the
root-product formula for the resultant, a difference of two real roots is
positive because the first root strip lies strictly to the right of the
second.  Nonreal roots can be grouped with their conjugates, making the
corresponding products positive squared absolute values.  Hence the
resultant itself is positive.
\end{proof}

It is useful to clear all dyadic denominators in the natural centered
coordinate.  For $h\geq1$, if $n_h=3(h-1)$, define
\begin{equation}\label{eq:meso-M-def}
 M_h(Y)=2^{n_h}N_h\left(\frac{Y-h-1}{2}\right)\in\mathbb Z[Y].
\end{equation}
Then
\[
 M_h(-Y)=(-1)^{h-1}M_h(Y),\qquad
 \operatorname {lc}M_h=\ell_h,
\]
and a linear change of variable gives the exact normalization
\begin{equation}\label{eq:meso-M-resultant}
 \operatorname {Res}_Y\bigl(M_d(Y),M_r(Y+d+r)\bigr)
 =2^{n_dn_r}\mathcal R_{d,r}.
\end{equation}

\begin{proposition}[Centered norm decomposition]\label{prop:meso-centered-norm}
For $d,r\geq2$, write
\[
 \widehat N_d(Y)=N_d\left(Y-\frac{d+1}{2}\right)
 =\begin{cases}
 A_d(Y^2),&d\ \text{odd},\\
 YA_d(Y^2),&d\ \text{even},
 \end{cases}
\]
where $A_d\in\mathbb Z[1/2][U]$, and set
\[
 G_{d,r}(Y)=N_r\left(Y+\frac{d-1}{2}\right),\qquad
 \Psi_{d,r}(Y^2)=G_{d,r}(Y)G_{d,r}(-Y).
\]
Then
\begin{equation}\label{eq:meso-centered-norm}
 |\mathcal R_{d,r}|=
 \begin{cases}
 |\operatorname {Res}_U(A_d,\Psi_{d,r})|,&d\ \text{odd},\\
 \left|N_r\bigl((d-1)/2\bigr)\right|
 |\operatorname {Res}_U(A_d,\Psi_{d,r})|,&d\ \text{even}.
 \end{cases}
\end{equation}
Moreover, the centered-coefficient formulas of
Lemma~\ref{lem:nv-centered-coefficients} give
\[
 A_d(U)=\ell_dU^m+u_dU^{m-1}+v_dU^{m-2}+\cdots
\]
with the evident adjustment when $d$ is even.
\end{proposition}

\begin{proof}
In the centered variable $Y=x+(d+1)/2$, the second polynomial is
$G_{d,r}(Y)$.  If $d$ is odd, the roots of $A_d(Y^2)$ occur in pairs
$\{\sqrt u,-\sqrt u\}$, and their two contributions to the root-product
formula multiply to $\Psi_{d,r}(u)$.  This is exactly the first
resultant in~\eqref{eq:meso-centered-norm}.  If $d$ is even, the
additional root $Y=0$ contributes $G_{d,r}(0)=N_r((d-1)/2)$.
The last assertion is the parity statement and the coefficient formulas
of Lemma~\ref{lem:nv-centered-coefficients}.
\end{proof}

Equivalently, in the finite $\mathbb Q$-algebra
\[
 K_d=\mathbb Q[X]/(N_d),\qquad \alpha=X\bmod N_d,
\]
one has
\begin{equation}\label{eq:meso-quotient-norm}
 \mathcal R_{d,r}
 =\ell_d^{3(r-1)}
 \operatorname {Nm}_{K_d/\mathbb Q}\bigl(N_r(\alpha+d)\bigr).
\end{equation}
No irreducibility or separability assumption is needed for this
determinantal norm identity.  For $d=2$,
$A_2(U)=34U+\tfrac32$; Proposition~\ref{prop:meso-centered-norm}
is precisely the rational linear factor times quadratic norm decomposition
recorded in Remark~\ref{rem:sep-res}.

\subsection{Half-integer center factors and their propagation}

The even centered factor gives a genuine divisibility lattice, rather than
only a numerical repetition pattern.

\begin{proposition}[Even-block center recurrence]\label{prop:meso-center-recurrence}
Fix an even integer $a\geq2$ and define
\[
 T_0^{(a)}=0,\qquad T_1^{(a)}=1,\qquad
 T_b^{(a)}=2^{3(b-1)}N_b\left(\frac{a-1}{2}\right)
 \quad(b\geq2).
\]
If $z=a+2b-1$, then
\begin{equation}\label{eq:meso-center-recurrence}
 T_{b+1}^{(a)}
 =\bigl(34z^3+102z^2+108z+40\bigr)T_b^{(a)}
 -z^6T_{b-1}^{(a)}.
\end{equation}
In particular, every $T_b^{(a)}$ is an integer.  Furthermore,
\begin{equation}\label{eq:meso-center-divides-first}
 T_b^{(a)}\mid\mathcal R_{a,b},
\end{equation}
and the odd part of $T_b^{(a)}$ divides $\mathcal R_{b,a+b}$.
By symmetry, if $b$ is even, the odd part of $T_a^{(b)}$ divides
both $\mathcal R_{a,b}$ and $\mathcal R_{a,a+b}$.
\end{proposition}

\begin{proof}
Substitute $x=(a-1)/2$ in the recurrence for $N_b$ and multiply by
$2^{3b}$.  Since $2(x+b)=a+2b-1=z$ and
\[
 8P(z/2)=34z^3+102z^2+108z+40,
\]
one obtains~\eqref{eq:meso-center-recurrence}.

Because $a$ is even, reflection shows that
$c=-(a+1)/2$ is a root of $N_a$.  Gauss's lemma gives
\[
 N_a(X)=(2X+a+1)Q_a(X),\qquad Q_a\in\mathbb Z[X].
\]
Resultant multiplicativity and the linear-root formula yield
\[
 \operatorname {Res}_X(2X+a+1,N_b(X+a))
 =2^{3(b-1)}N_b(c+a)=T_b^{(a)},
\]
which proves~\eqref{eq:meso-center-divides-first}.

Let $p^e$ be any odd prime power dividing $T_b^{(a)}$.  Since $2$ is a
unit modulo $p^e$, this says
\[
 N_b(t)=0\pmod {p^e},\qquad t=(a-1)/2.
\]
The restart identity is valid without division in the present situation:
the two sides satisfy the same recurrence and hence
\[
 N_{a+b}(c)=N_{a+1}(c)N_b(c+a)
             =N_{a+1}(c)N_b(t)=0\pmod {p^e}.
\]
Put $y=-t-b-1$.  Reflection gives
$N_b(y)=(-1)^{b-1}N_b(t)=0$, while $y+b=c$, so
$N_{a+b}(y+b)=0$.  The Sylvester--B\'ezout identity for the resultant,
evaluated at $y\in\mathbb Z/p^e\mathbb Z$, gives
$p^e\mid\mathcal R_{b,a+b}$.  Taking all odd prime powers proves the
claim.  The symmetric assertion follows from
$\mathcal R_{a,b}=\mathcal R_{b,a}$.
\end{proof}

There is also a simultaneous center factor for every pair.  With
$T_b^{(a)}$ as above, define
\[
 \mathfrak C_{d,r}=
 \begin{cases}
 1,&d,r\ \text{odd},\\
 |T_r^{(d)}|,&d\ \text{even},\ r\ \text{odd},\\
 |T_d^{(r)}|,&d\ \text{odd},\ r\ \text{even},\\
 \displaystyle\frac{|T_r^{(d)}T_d^{(r)}|}{2(d+r)},
     &d,r\ \text{even}.
 \end{cases}
\]
Then $\mathfrak C_{d,r}$ is an integer and
\begin{equation}\label{eq:meso-center-factor}
 \mathfrak C_{d,r}\mid\mathcal R_{d,r}.
\end{equation}
Indeed, factor every available central linear factor before applying
resultant multiplicativity.  When both indices are even, the two linear
factors have resultant $2(d+r)$, which accounts for the denominator and
prevents counting their mutual resultant twice.

\subsection{The diagonal square law}

For the diagonal family, the centered pairing becomes an exact square.

\begin{theorem}[Diagonal square law]\label{thm:meso-diagonal-square}
For $d\geq2$, let $n=3(d-1)$ and put
\[
 f_d(T)=N_d\left(T-\frac{2d+1}{2}\right)
       =E_d(T^2)+TO_d(T^2),
 \qquad E_d,O_d\in\mathbb Q[U].
\]
Then
\begin{equation}\label{eq:meso-diagonal-square-rational}
 \mathcal R_{d,d}
 =(-1)^{d-1}2^n\ell_d
 N_d\left(-\frac{2d+1}{2}\right)
 \operatorname {Res}_U(O_d,E_d)^2.
\end{equation}
Moreover,
\[
 D_d=2^nN_d(-1/2)\in\mathbb Z,
 \qquad Q_d=\operatorname {Res}_U(O_d,E_d)\in\mathbb Z,
\]
and therefore
\begin{equation}\label{eq:meso-diagonal-square-integral}
 \boxed{\mathcal R_{d,d}=\ell_d|D_d|Q_d^2.}
\end{equation}
For every odd prime $p$,
\begin{equation}\label{eq:meso-diagonal-valuation}
 v_p(\mathcal R_{d,d})
 =v_p(\ell_d)+v_p(D_d)+2v_p(Q_d).
\end{equation}
In particular, if $p\nmid\ell_dD_d$, then
$v_p(\mathcal R_{d,d})$ is even.
\end{theorem}

\begin{proof}
We first record a general identity.  Let
\[
 f(T)=aT^n+\cdots=E(T^2)+TO(T^2),
\]
and suppose the two parity parts have the full degrees forced by nonzero
coefficients in degrees $n$ and $n-1$.  If $\alpha_i$ are the roots of
$f$, then
\[
 f(-\alpha_i)=-2\alpha_iO(\alpha_i^2),
\]
and the root-product formula gives
\begin{equation}\label{eq:meso-square-step1}
 \operatorname {Res}_T(f(T),f(-T))
 =2^na^{n-1}f(0)\prod_iO(\alpha_i^2).
\end{equation}
If $s=\deg O$, $e=\deg E$, and $b=\operatorname {lc}O$, then
\begin{align*}
 a^{2s}\prod_iO(\alpha_i^2)
 &=\operatorname {Res}_T(f(T),O(T^2))\\
 &=b^{n-2e}\operatorname {Res}_U(O,E)^2.
\end{align*}
For $n=2m$, one has $(e,s)=(m,m-1)$.  For $n=2m+1$, one has
$(e,s)=(m,m)$ and $b=a$.  In either case, substitution into
\eqref{eq:meso-square-step1} gives
\begin{equation}\label{eq:meso-general-square}
 \operatorname {Res}_T(f(T),f(-T))
 =2^na f(0)\operatorname {Res}_U(O,E)^2.
\end{equation}

For $f=f_d$, the natural centered parity of $N_d$ shows that the
coefficient of $T^{n-1}$ is $-nd\ell_d/2$, so the stated degree
conditions hold.  Reflection gives
\[
 N_d(T-1/2)=(-1)^{d-1}f_d(-T).
\]
After the translation $T=x+(2d+1)/2$, scaling the second argument of
the resultant contributes
$(-1)^{(d-1)n}=(-1)^{d-1}$.  Equation
\eqref{eq:meso-general-square} is therefore exactly
\eqref{eq:meso-diagonal-square-rational}.

It remains only to justify the integral form.  Integrality of $D_d$ is
immediate from evaluating an integer polynomial of degree $n$ at a
half-integer.  Put $c=2d+1$ and $V=X(X+c)$.  Reduction by the monic
relation $X^2=V-cX$ gives the free-module decomposition
\[
 \mathbb Z[X]=\mathbb Z[V]\oplus X\mathbb Z[V].
\]
Thus there are $A,B\in\mathbb Z[V]$ such that
\[
 N_d(X)=A(V)+XB(V).
\]
Since $T=X+c/2$ and $U=T^2$, one has $V=U-c^2/4$, and comparison with
$f_d(T)=E_d(U)+TO_d(U)$ gives
\[
 O_d(U)=B(U-c^2/4),
 \qquad
 E_d(U)=A(U-c^2/4)-\frac c2B(U-c^2/4).
\]
Translation invariance reduces $Q_d$ to
$\operatorname {Res}_V(B,A-\tfrac c2B)$.  In fact $\deg A=\deg E_d$.
When $n=2m$, this follows from $\deg E_d=m$ and $\deg B=m-1$.
When $n=2m+1$, the leading coefficients computed above give
\[
 \operatorname {lc}B=\ell_d,\qquad
 \operatorname {lc}A
 =\frac{c-nd}{2}\ell_d,
\]
and $c-nd=-3d^2+5d+1\ne0$ for every integer $d\geq2$.  The
root-product formula therefore gives
\[
 Q_d=\operatorname {Res}_U(O_d,E_d)
 =\operatorname {Res}_V(B,A-\tfrac c2B)
 =\operatorname {Res}_V(B,A)\in\mathbb Z.
\]
Finally, reflection gives
\[
 \left|N_d\left(-\frac{2d+1}{2}\right)\right|=|N_d(-1/2)|.
\]
Proposition~\ref{prop:meso-root-strip} supplies the sign, and
\eqref{eq:meso-diagonal-square-integral} follows.  The valuation formula
is immediate.
\end{proof}

In odd characteristic the affine diagonal roots can also be read off
exactly.  In the coordinate used above,
\[
 f_d(T)=f_d(-T)=0
 \quad\Longleftrightarrow\quad
 E_d(T^2)=0,\quad TO_d(T^2)=0.
\]
Thus the number of distinct $\mathbb F_p$-roots is
\begin{equation}\label{eq:meso-diagonal-root-count}
 \mathbf1_{p\mid f_d(0)}
 +2\#\{u\in(\mathbb F_p^\times)^2:E_d(u)=O_d(u)=0\}.
\end{equation}
This distinguishes an actual split collision from a nonsplit quadratic
factor of the resultant.

\subsection{What resultant divisibility does and does not mean}

\begin{proposition}[Projective validity domain]\label{prop:meso-resultant-validity}
For integer polynomials of their nominal degrees,
$p\mid\mathcal R_{d,r}$ if and only if their nominal homogenizations
have a common point in $\mathbb P^1(\overline{\mathbb F}_p)$.  If the
two leading coefficients do not both vanish modulo $p$, this is
equivalent to an affine common root in $\overline{\mathbb F}_p$.
It is not, in general, equivalent to a common root in $\mathbb F_p$.
The universally valid implication is only
\[
 \{x\in\mathbb F_p:N_d(x)=N_r(x+d)=0\}\ne\varnothing
 \quad\Longrightarrow\quad p\mid\mathcal R_{d,r}.
\]
\end{proposition}

Two mesoscopic examples show both failures of the converse.
For $d=r=2$ and $p=17$,
\[
 N_2(X)\equiv5(2X+3),\qquad
 N_2(X+2)\equiv5(2X+7)\pmod {17}.
\]
They have no affine common root, although
$17^3\mid\mathcal R_{2,2}$; both nominal cubics lost their leading
terms and meet at infinity.  On the other hand, for $d=r=3$ and
$p=110629$, the leading coefficient is a unit and
\[
 \gcd\bigl(N_3(X),N_3(X+3)\bigr)
 =X^2+7X-15495.
\]
Its discriminant is $62029$, a nonsquare modulo $110629$.  Hence there
is an affine common root over $\overline{\mathbb F}_p$ but none over
$\mathbb F_p$, even though $110629^2\mid\mathcal R_{3,3}$.

The centered coefficients identify the complete lattice of automatic
infinity edges.  Recall
\[
 \ell_0=0,\quad\ell_1=1,\quad
 \ell_{n+1}=34\ell_n-\ell_{n-1}.
\]
For $p\geq5$, choose $\lambda\in\mathbb F_{p^2}$ with
$\lambda+\lambda^{-1}=34$.  Then
\[
 \ell_n=\frac{\lambda^n-\lambda^{-n}}
              {\lambda-\lambda^{-1}}.
\]
If $\rho_p=\operatorname {ord}(\lambda^2)$, then
\begin{equation}\label{eq:meso-rank-lattice}
 p\mid\ell_n\quad\Longleftrightarrow\quad\rho_p\mid n,
 \qquad \rho_p\mid p-\left(\frac2p\right).
\end{equation}
Consequently every pair with $\rho_p\mid d,r$ is in the raw resultant
support, independently of affine collisions.  If
$M=\lfloor H/\rho_p\rfloor$, these pairs alone number
\begin{equation}\label{eq:meso-infinity-count}
 V_p^{\mathrm{infinity}}(H)=\frac{M(M-1)}2.
\end{equation}

There is a legitimate centered-coefficient transfer here, but it detects
only the false edge at infinity.  Suppose $p\geq7$ and $d<p$.  If
$\ell_d=0$, Cassini gives $\ell_{d-1}^2=1$, while
Lemma~\ref{lem:nv-centered-coefficients} gives
\[
 u_d=\frac{5d\ell_{d-1}}{256}\ne0\pmod p.
\]
Because the natural centered polynomial contains only every other power,
its degree drops by exactly two.  If $\rho_p\mid d,r$, the two nominal
homogenizations have a common factor of order two at infinity; equivalently,
the nominal Sylvester matrix has nullity at least two modulo $p$.  Its Smith
normal form therefore gives
\begin{equation}\label{eq:meso-infinity-square}
 p^2\mid\mathcal R_{d,r}.
\end{equation}

For example, $p=665857$ has $\rho_p=8$ and
$H=\lfloor\sqrt p\rfloor=816$.  Formula
\eqref{eq:meso-infinity-count} already gives $5151$ raw resultant pairs,
without producing a single affine collision.  Thus the unsaturated integer
resultant support is not the correct collision graph.

The centered/Pell method from Theorem~\ref{thm:nv-range} cannot be used
pointwise in the opposite direction.  At $(d,r,p)=(2,2,71)$ and
$x=-5/2$, the two natural centered variables are $-1$ and $1$, and
\[
 \widehat N_2(Y)=34Y^3+\frac32Y,\qquad
 \widehat N_2(\pm1)=\pm\frac{71}{2}\equiv0\pmod {71},
\]
while $\ell_2=34\ne0\pmod {71}$.  The proof of
Theorem~\ref{thm:nv-range} compares coefficients only after assuming
that an entire certificate polynomial is identically zero.  Two point-value
equations do not make its coefficients $L,C_1,C_2,C_4$ vanish.  In
particular, its $C_4$ formula is valid only inside the stated danger case
and cannot be transferred to this collision.

\subsection{The two missing counting lemmas}

For an integer $H\geq2$, define
\[
 \mathscr D_H=\{(d,r)\in\mathbb Z_{\geq2}^2:d+r\leq H\},
\]
\[
 X_{d,r}(p)=\{x\in\mathbb F_p:
 N_d(x)=N_r(x+d)=0\},\qquad m_{d,r}(p)=|X_{d,r}(p)|.
\]
The collision energy is the weighted count
\[
 \mathcal E_p(H)=\sum_{(d,r)\in\mathscr D_H}m_{d,r}(p).
\]
Define also the simple affine support
\[
 W_p(H)=\#\{(d,r)\in\mathscr D_H:m_{d,r}(p)>0\}.
\]
The raw resultant support
\[
 V_p(H)=\#\{(d,r)\in\mathscr D_H:p\mid\mathcal R_{d,r}\}
\]
is the edge count of a still larger projective-support graph.  Thus
\[
 W_p(H)\leq\mathcal E_p(H),\qquad W_p(H)\leq V_p(H).
\]
Remark~\ref{rem:collision} identified $\mathcal E_p(H)$ with the edge
count $W_p(H)$ by forgetting that one pair can have several root witnesses.
For example, at $p=3109$ the pair $(d,r)=(8,8)$ has two distinct
$\mathbb F_p$ common roots, so it contributes one to $W_p$ and two to
$\mathcal E_p$.  The examples above exhibit the separate failure of
$W_p(H)=V_p(H)$.

The same distinction changes the averaging statement in that remark.
The Hadamard argument controls only
\[
 \sum_{X<p\leq2X}V_p(H)\ll\frac{H^4\log H}{\log X}.
\]
After inserting the available multiplicity bound, its valid consequence is
\[
 \sum_{X<p\leq2X}\mathcal E_p(H)
 \ll\frac{H^5\log H}{\log X}.
\]
Since there are $\asymp X/\log X$ primes in the interval, this proves
$\mathcal E_p(H)\ll H^{1/2}$ for all but $o(X/\log X)$ of them only in
the range $H\leq X^{2/9-\varepsilon}$, rather than the stated
$X^{2/7-\varepsilon}$ range.  The latter exponent applies to simple
support, not to the collision energy.

The existing gcd estimate gives
\[
 m_{d,r}(p)\leq3(\min(d,r)-1).
\]
Accordingly, define the degree-weighted support
\[
 \mathcal W_p(H)=
 \sum_{(d,r)\in\mathscr D_H}
 (\min(d,r)-1)\mathbf1_{p\mid\mathcal R_{d,r}}.
\]
Then the rigorous implication is
\begin{equation}\label{eq:meso-weighted-energy}
 \mathcal E_p(H)\leq3\mathcal W_p(H).
\end{equation}
In particular,
\begin{equation}\label{eq:meso-V-only}
 V_p(H)\ll H\quad\Longrightarrow\quad
 \mathcal E_p(H)\ll H^2,
\end{equation}
not $H^{3/2}$.  This loss is real at the level of the proposed graph
argument: take $\asymp H$ exceptional pairs with $d,r\asymp H$, make them
a matching (maximum degree one and codegree zero), and assign
$\asymp H$ distinct root witnesses to every edge.  This obeys the available
per-pair degree bound and has energy $\asymp H^2$.  The ordinary
K\H{o}v\'ari--S\'os--Tur\'an theorem applies to a simple graph and cannot
control those parallel root witnesses.

Thus a sufficient replacement for G1 is the genuinely weighted estimate
\begin{equation}\label{eq:meso-missing-weighted}
 \boxed{\mathcal W_p(H)\ll H^{3/2}.}
\end{equation}
Alternatively, together with $V_p(H)\ll H$, it would suffice to prove an
average actual gcd-degree bound $O(H^{1/2})$ on the exceptional pairs, or
a corresponding multiplicity-tail estimate.

There is a second, independent gap between energy and the first moment
$R_p(H)=\sum_{h=2}^H r_p(h)$.  Let
\[
 k_x=\#\{2\leq h\leq H:N_h(x)=0\},\qquad
 \partial_H=\{-2,-3,\ldots,-H\}\subseteq\mathbb F_p.
\]
For $x\notin\partial_H$, Proposition~\ref{prop:bezout} maps every pair
$2\leq d<e\leq H$ counted by $\binom{k_x}{2}$ injectively to the energy
term $(x,d,e-d)$.  Indeed $e-d=1$ cannot occur: in that case the B\'ezout
right-hand side is $S_d(x)^6N_1=S_d(x)^6\ne0$.  Thus every resulting
pair has $e-d\geq2$ and lies in $\mathscr D_H$.  Hence
\begin{equation}\label{eq:meso-column-pairs}
 \sum_{x\notin\partial_H}\binom{k_x}{2}\leq\mathcal E_p(H).
\end{equation}
The exclusion of $\partial_H$ is essential.  The unsaturated degree claim
following Proposition~\ref{prop:bezout} fails at cut-edge roots: for
$(d,e,p)=(3,4,73)$,
\[
 \gcd_{\mathbb F_{73}[X]}(N_3,N_4)=X+3,
\]
whereas $N_{e-d}=N_1=1$.  Thus that claim is valid only for the saturated
gcd, or as an away-from-boundary root count; the use above has precisely
that domain.

The same example invalidates the proof of the unrestricted column bound in
Proposition~\ref{prop:column}.  Since $b_2=73$, the endpoint formula gives
\[
 N_h(-3)\equiv0\pmod {73}\qquad(h\geq3).
\]
At the singular transition the recurrence has reached two consecutive
zeros, so the entire later tail remains zero; it cannot be charged as an
$O(1)$ transition cost.  Thus that proof does not establish the proposition
or the rectangular incidence bound derived from it; a valid version needs
an explicit boundary term or an unpolluted nonsingular initial state.  The
present argument excludes such columns from
\eqref{eq:meso-column-pairs} and retains their full contribution in
$R_p^\partial(H)$.
But columns with $k_x=1$ contribute nothing to either side of
\eqref{eq:meso-column-pairs}.  Energy alone cannot bound their total mass.

For a precise sufficient condition, set $K=\lceil\sqrt H\rceil$ and
\[
 R_p^\partial(H)=\sum_{x\in\partial_H}k_x,
 \qquad
 L_p(H)=\sum_{\substack{x\notin\partial_H\\0<k_x<K}}k_x.
\]
Since $k\leq2\binom{k}{2}/(K-1)$ for $k\geq K$,
\eqref{eq:meso-column-pairs} gives the deterministic inequality
\begin{equation}\label{eq:meso-amplification-exact}
 \boxed{
 R_p(H)\leq R_p^\partial(H)+L_p(H)
 +\frac{2}{\lceil\sqrt H\rceil-1}\mathcal E_p(H).}
\end{equation}
Thus the missing low-fiber amplification is
\begin{equation}\label{eq:meso-missing-amplification}
 R_p^\partial(H)+L_p(H)\ll H,
\end{equation}
or, in compressed form,
\[
 R_p(H)\ll H+H^{-1/2}\mathcal E_p(H).
\]
This is a lower-participation assertion; K\H{o}v\'ari--S\'os--Tur\'an is
an upper-bound theorem and does not supply it.  The boundary term cannot
be discarded, since the endpoint factorization of $N_h(-j)$ can make a
cut-edge column large when $p\mid b_{j-1}$.

Combining the two genuinely missing estimates
\eqref{eq:meso-missing-weighted} and
\eqref{eq:meso-missing-amplification} would give
\[
 \mathcal E_p(H)\ll H^{3/2},\qquad R_p(H)\ll H.
\]
Only then does the valid last step
\[
 Z(p)\leq\left\lceil\frac pH\right\rceil+R_p(H-1)
\]
give $Z(p)\ll\sqrt p$ at $H=\lfloor\sqrt p\rfloor$.

For completeness, suppose only the proposed G2 support estimate
\[
 V_p(H)\ll H^{2-\delta}
\]
is available.  The present ingredients give merely
\[
 \mathcal E_p(H)\ll H^{3-\delta}.
\]
Partition $\{0,\ldots,p-1\}$ into nonwrapping consecutive blocks of
length at most $H$, and let their Ap\'ery zero counts be $z_i$.  Three
zeros $m<m+d<m+d+r$ in one block have $d,r\geq2$ by
Lemma~\ref{lem:no-consec}.  All intervening recurrence denominators are
units modulo $p$, so the gap-polynomial identity gives the distinct energy
witness $(m,d,r)$.  Hence
\[
 \sum_i\binom{z_i}{3}\leq\mathcal E_p(H)
\]
and H\"older's inequality yields
\begin{equation}\label{eq:meso-triple-endpoint}
 Z(p)\ll\frac pH+p^{2/3}H^{(1-\delta)/3}.
\end{equation}
If the G2 estimate holds uniformly for $H\leq\sqrt p$, then, after
combining with the existing $p^{2/3}$ bound,
\[
 Z(p)\ll
 \begin{cases}
 p^{2/3},&0<\delta\leq1,\\
 p^{(5-\delta)/6},&1<\delta\leq2.
 \end{cases}
\]
Thus even $V_p(H)\ll H$ ($\delta=1$) recovers only the existing exponent
from the currently proved inputs.  Even a separately proved
$\mathcal E_p(H)\ll H^{3/2}$ would give only $Z(p)\ll p^{7/12}$ by this
triple-counting route unless the low-fiber amplification is also proved.

Finally, the recurrence does not close on the scalar resultants.  In the
algebra $K_d$ of~\eqref{eq:meso-quotient-norm}, put
$u_r=N_r(\alpha+d)$.  Then
\[
 u_{r+1}=P(\alpha+d+r)u_r-(\alpha+d+r)^6u_{r-1}.
\]
The norm of a sum depends on all polarized mixed norms, so
$\operatorname {Nm}(u_r)$ and $\operatorname {Nm}(u_{r-1})$ do not
determine $\operatorname {Nm}(u_{r+1})$.  Modulo $p$, different values of
$r$ may also vanish in different residue-field components, whereas restart
requires the same root witness.  This root-witness alignment, the weighted
support estimate~\eqref{eq:meso-missing-weighted}, and the low-fiber
amplification~\eqref{eq:meso-missing-amplification} are the exact missing
lemmas left by the present attack.
